% =====================================================================
% Dataset, preprocessing and baseline CNN  (author: Abhishek Singh, Stage 1)
% Every number below comes from results/ or data/splits/ in the repository:
%   results/stage1_eda/*.csv, data/splits/split_summary.json,
%   results/stage1_baseline/settings.json, results/stage1_metrics.csv
% Figures to upload to Overleaf (keep these paths):
%   figures/stage1/sample_grid.png          <- results/stage1_eda/sample_grid.png
%   figures/stage1/image_sizes.png          <- results/stage1_eda/image_sizes.png
%   figures/stage1/training_curves.png      <- results/stage1_baseline/training_curves.png
%   figures/stage1/confusion_matrices.png   <- results/stage1_baseline/confusion_matrices.png
% =====================================================================
\section{Dataset and Preprocessing}
\label{sec:data}

\subsection{Dataset}
\label{sec:data-dataset}
We use the \emph{Chest X-Ray Images (Pneumonia)} dataset~\cite{mooney2018kaggle}, the Kaggle release of the paediatric chest X-rays
collected by Kermany \etal~\cite{kermany2018identifying} at the Guangzhou Women and Children's Medical Center.
It contains 5,856 frontal chest X-rays in two classes: 1,583 \textsc{normal} and 4,273 \textsc{pneumonia}
(2,780 bacterial and 1,493 viral, which we treat as one class).
The images are split into official \texttt{train}, \texttt{val} and \texttt{test} folders (\cref{tab:data-splits}).
The download also contains a duplicate copy of all images in a nested \texttt{chest\_xray/} folder and a \texttt{\_\_MACOSX/} folder; we read only the top-level folders.
\Cref{fig:samples} shows example images.

\begin{figure*}[t]
  \centering
  \includegraphics[width=\linewidth]{figures/stage1/sample_grid.png}
  \caption{Random training images (fixed seed) with their original size. Text markers such as the letter \emph{R} appear in both classes.
           Pneumonia images are often wider relative to their height (\cref{sec:data-eda}).}
  \label{fig:samples}
\end{figure*}

\paragraph{Patients.}
Several images can belong to the same child (up to 30 for one patient), so images of one patient must never be spread across training and evaluation sets.
The filenames encode a patient identifier, from which we derive a patient ID:
\texttt{person\{N\}\_\{bacteria|virus\}} for pneumonia images and \texttt{IM-\{XXXX\}} or \texttt{NORMAL2-IM-\{XXXX\}} for normal images.
With this rule the official folders share no patients, and the dataset contains 4,097 patients (1,444 normal, 2,653 pneumonia).
If only \texttt{person\{N\}} were used, 170 identifiers would appear in both the official training and test folders.
Including the infection type removes all of this overlap, which strongly suggests that \texttt{person\{N\}\_bacteria} and \texttt{person\{N\}\_virus} are different children.
This is an inference from the file names and not a fact documented by the dataset authors (\cref{sec:data-limits}).

\subsection{Data splits}
\label{sec:data-splits}
The official validation folder contains only 16 images, which is too few for model selection.
We keep the official test folder unchanged as the test set and build a new validation set from the remaining 5,232 images
(official \texttt{train} plus \texttt{val}).
The split is \emph{patient-level} and stratified by class: for each class, the patients are shuffled with a fixed seed (42)
and assigned to validation until it holds 12.5\% of that class's images; all images of a patient go to the same set.
As a safety margin, the training/validation split groups pneumonia images by \texttt{person\{N\}} alone, keeping a person's bacterial and viral images on the same side,
so that the two sets share no identifier even if both types belonged to the same child.
The resulting sets (\cref{tab:data-splits}) share no patients (and training and validation share no \texttt{person\{N\}}), and both training and validation contain 74.2\% pneumonia images.
We first used scikit-learn's \texttt{StratifiedGroupKFold}, but with the same seed it produced different splits under different library versions.
We therefore implemented the split directly with NumPy's legacy \texttt{RandomState}, whose random stream is fixed, so the split is identical on every platform.
The split is stored as CSV files in the repository, and every stage of the project uses these exact files.

\begin{table}[t]
  \centering
  \small
  \begin{tabular}{lrrrr}
    \toprule
    Set & Images & Normal & Pneumonia & Patients \\
    \midrule
    \multicolumn{5}{l}{\emph{Official folders}} \\
    \texttt{train} & 5,216 & 1,341 & 3,875 & 3,654 \\
    \texttt{val}   &    16 &     8 &     8 &    15 \\
    \texttt{test}  &   624 &   234 &   390 &   428 \\
    \midrule
    \multicolumn{5}{l}{\emph{Our splits (used in all experiments)}} \\
    Train      & 4,578 & 1,180 & 3,398 (74.2\%) & 3,218 \\
    Validation &   654 &   169 &   485 (74.2\%) &   451 \\
    Test       &   624 &   234 &   390 (62.5\%) &   428 \\
    \bottomrule
  \end{tabular}
  \caption{Images and patients per set. Our validation set is a patient-level, stratified 12.5\% of the official \texttt{train} and \texttt{val} folders; the official test folder is unchanged.}
  \label{tab:data-splits}
\end{table}

\subsection{Exploratory analysis and potential shortcuts}
\label{sec:data-eda}
The images vary widely in size: widths range from 384 to 2,916 pixels and heights from 127 to 2,713 pixels.
Two properties of the raw files differ systematically between the classes (\cref{fig:sizes}):
\begin{itemize}
  \item \textbf{Size and aspect ratio.} Normal images are larger (median $1654\times1323$ pixels) than pneumonia images (median $1160\times776$),
        and their median width-to-height ratio is about 1.2, compared with about 1.5 for pneumonia images.
        A plain resize to a square input therefore squeezes the two classes by different amounts, so a model could learn the distortion instead of the lung findings.
  \item \textbf{Colour mode.} All 283 RGB files in the dataset are pneumonia images; every normal image is stored as single-channel grayscale.
\end{itemize}
Both properties are potential shortcuts that exist before any model is trained.
We remove the colour-mode cue by converting every image to grayscale (\cref{sec:data-preproc}).
We keep the plain resize: padding to a square would encode the aspect ratio in the width of the black bars and add artificial borders,
and cropping could cut off the edges of the lungs.
We therefore treat the aspect ratio as a known confound when interpreting the results.

\begin{figure}[t]
  \centering
  \includegraphics[width=\linewidth]{figures/stage1/image_sizes.png}
  \caption{Image size (left) and aspect ratio (right) per class. Pneumonia images are smaller and relatively wider, so a square resize distorts the two classes differently.}
  \label{fig:sizes}
\end{figure}

\subsection{Preprocessing and augmentation}
\label{sec:data-preproc}
All models in this project share the same preprocessing.
Each image is converted to grayscale, copied to three channels (so that ImageNet-pretrained networks can be used in later stages),
resized to $224\times224$ pixels without padding or cropping, and normalised with the ImageNet mean and standard deviation.
During training only, images are rotated by a random angle of up to $\pm10^\circ$.
We deliberately do not use horizontal flips: flipping a chest X-ray moves the heart to the other side and mirrors the laterality markers,
which are among the shortcuts this project studies.
The classes are imbalanced (74.2\% pneumonia in the training set).
We compensate in the loss with a positive-class weight equal to the ratio of normal to pneumonia images in the training set, $w_{+} = 1180/3398 = 0.347$.
Data loading, shuffling and augmentation are seeded, so training runs are reproducible.

\section{Baseline CNN}
\label{sec:baseline}

\paragraph{Architecture.}
As a baseline we train a custom CNN from scratch.
It has five convolutional blocks with 32, 64, 128, 256 and 256 channels.
Each block applies two $3\times3$ convolutions, each followed by batch normalisation~\cite{ioffe2015batch} and a ReLU, and then $2\times2$ max pooling.
The blocks are followed by global average pooling, dropout ($p=0.3$) and a single linear output: the logit of the pneumonia class.
The network has 2.35\,M parameters.
The last convolutional layer produces a $14\times14$ feature map, which Grad-CAM uses to produce the heatmaps in Stage~3.

\paragraph{Training.}
The model is trained with binary cross-entropy on the logit, using the positive-class weight from \cref{sec:data-preproc},
and the Adam optimiser~\cite{kingma2015adam} (learning rate $10^{-3}$, weight decay $10^{-4}$, batch size 32).
The learning rate is halved when the validation loss has not improved for two epochs.
Training stops when the validation loss has not improved for five epochs (at most 30 epochs), and we keep the weights with the lowest validation loss.
Model selection uses only the validation set: the test set is evaluated once, with the final model and a fixed decision threshold of 0.5.
The best validation loss was reached at epoch 3, and training stopped after epoch 8 (\cref{fig:curves}).
Training took about two minutes on a single NVIDIA T4 GPU with mixed precision, using PyTorch~\cite{paszke2019pytorch}.

\begin{figure}[t]
  \centering
  \includegraphics[width=\linewidth]{figures/stage1/training_curves.png}
  \caption{Baseline training: loss (left) and validation AUC and F1 (right) per epoch. The weights from epoch 3 (lowest validation loss) are kept.}
  \label{fig:curves}
\end{figure}

\paragraph{Results.}
\Cref{tab:baseline} reports the baseline on the validation and test sets, with pneumonia as the positive class.
On validation, the model reaches 96.8\% accuracy and an AUC of 0.992.
On the test set, accuracy drops to 80.1\% (AUC 0.902).
The confusion matrices (\cref{fig:cm}) show where the difference comes from.
The model misses only 7 of the 390 pneumonia test images (recall 0.982), but it labels 117 of the 234 normal test images as pneumonia,
so it correctly identifies only 50.0\% of normal test images, compared with 95.9\% on validation.
Our validation set is drawn from the same pool as the training data, while the official test set appears to differ in distribution,
and the model's errors on it are almost entirely false alarms on normal images.
We did not tune the threshold or any other setting on the test set.
DenseNet-121 later raises test accuracy to 0.873 and AUC to 0.971, and reduces false positives on normal images from 117 to 78 (\cref{sec:densenet-results}).
This gap is a starting point for the attention analysis: if the false positives are driven by regions outside the lungs,
the Grad-CAM metric and the masking experiment should reveal it.

\begin{table}[t]
  \centering
  \small
  \begin{tabular}{lccccc}
    \toprule
    Set & Acc. & Prec. & Recall & F1 & AUC \\
    \midrule
    Validation & 0.968 & 0.985 & 0.971 & 0.978 & 0.992 \\
    Test       & 0.801 & 0.766 & 0.982 & 0.861 & 0.902 \\
    \bottomrule
  \end{tabular}
  \caption{Custom CNN baseline (threshold 0.5, pneumonia = positive class). Source: \texttt{results/stage1\_metrics.csv}.}
  \label{tab:baseline}
\end{table}

\begin{figure}[t]
  \centering
  \includegraphics[width=\linewidth]{figures/stage1/confusion_matrices.png}
  \caption{Confusion matrices of the baseline on validation (left) and test (right). On the test set, 117 of 234 normal images are predicted as pneumonia.}
  \label{fig:cm}
\end{figure}

\subsection{Limitations of the data}
\label{sec:data-limits}
% Can be moved into the shared Limitations / Discussion section.
The patient IDs are inferred from file names (\cref{sec:data-dataset}); zero overlap between the official folders is strong evidence, but not proof, that our split separates patients.
The aspect-ratio difference between the classes (\cref{sec:data-eda}) is not removed by our preprocessing, and lung masking (Stage~4) does not remove it either, because the masked lungs are still distorted.
Finally, all images come from a single medical centre, so none of our test results measure how well the models generalise to other hospitals.
